% chapters/05/praktikum.tex
\begin{praktikumbox}
Praktikum ini mengikuti Tugas Praktikum 03 pada RPS: menggerakkan karakter
dengan tombol WASD atau tombol panah pada kecepatan tetap (piksel per detik),
lalu membandingkan perilakunya pada 30 dan 60 bingkai per detik. Isilah lembar
kerja pada Tabel~\ref{tab:lembar-praktikum-03}.

\begin{langkah}
  \item \textbf{Jalankan program gerak.} Ketik ulang \code{ch05_gerak.py} dari
        Daftar~\ref{lst:ch05-gerak} ke dalam folder \code{code}. Jalankan, lalu
        gerakkan persegi dengan tombol WASD dan dengan tombol panah. Pastikan
        keduanya bekerja.

  \item \textbf{Amati penyelarasan posisi dan \code{Rect}.} Perhatikan dua
        baris pertama pada jendela. Gerakkan objek perlahan, lalu perhatikan
        bahwa nilai pecahan berubah dengan halus sedangkan nilai bulat melompat
        satu piksel setiap kali.

  \item \textbf{Ukur kecepatan pada 30 bingkai per detik.} Jalankan program
        dengan perintah berikut, lalu catat kecepatan terukurnya.
\begin{kodekeluaran}
python ch05_gerak.py --fps 30 --frames 18 --auto-direction kanan
\end{kodekeluaran}

  \item \textbf{Ukur kecepatan pada 60 dan 120 bingkai per detik.} Ulangi
        langkah sebelumnya dengan \code{--fps 60 --frames 36} dan
        \code{--fps 120 --frames 72}, yaitu jumlah bingkai yang menghasilkan
        lama berjalan kurang lebih sama. Bandingkan kecepatan terukurnya.

  \item \textbf{Uji gerak diagonal.} Jalankan program dengan
        \code{--auto-direction diagonal}, lalu bandingkan kecepatan terukurnya
        dengan gerak lurus. Sesudah itu, hapus pemanggilan \code{normalize()}
        pada fungsi \code{update_position}, jalankan kembali, dan catat
        perbedaannya. Kembalikan baris itu sesudah selesai.

  \item \textbf{Tambahkan objek kedua.} Tambahkan satu persegi lagi yang
        bergerak dengan kecepatan berbeda, misalnya 120 piksel per detik
        dengan tombol IJKL. Pastikan kedua objek bergerak dengan kecepatan
        masing-masing tanpa saling mempengaruhi.

  \item \textbf{Tulis catatan praktikum} yang memuat tabel pengukuran Anda
        beserta kesimpulan mengenai hubungan antara laju bingkai dan kecepatan
        objek.

  \item \textbf{Simpan pekerjaan} ke repositori Git dengan pesan penyimpanan
        yang menerangkan isinya.
\end{langkah}
\end{praktikumbox}

\begin{table}[htbp]
\centering
\caption{Lembar kerja Praktikum 03}
\label{tab:lembar-praktikum-03}
{\setlength{\tabcolsep}{4pt}\footnotesize%
\begin{tabularx}{\textwidth}{@{}C{2.6cm} C{2.6cm} C{2.8cm} Y@{}}
\toprule
\headrow \textbf{Batas laju} & \textbf{Waktu berjalan} & \textbf{Jarak tempuh} & \textbf{Kecepatan terukur} \\
\midrule
30 FPS, 18 bingkai & & & \\[1.4em]
\addlinespace
60 FPS, 36 bingkai & & & \\[1.4em]
\addlinespace
120 FPS, 72 bingkai & & & \\[1.4em]
\addlinespace
\multicolumn{4}{@{}p{\textwidth}@{}}{\footnotesize\itshape
Kecepatan yang ditetapkan: 200 piksel per detik, arah lurus.} \\
\addlinespace
\multicolumn{4}{@{}p{\textwidth}@{}}{\footnotesize\itshape
Arah diagonal tanpa \code{normalize()}: kecepatan terukur \dotfill} \\[1.4em]
\bottomrule
\end{tabularx}}
\end{table}

\begin{catatabox}
Apabila kecepatan terukur Anda berbeda-beda secara berarti pada ketiga laju
bingkai, periksa lebih dahulu apakah \code{delta_seconds} benar-benar dipakai
sebagai pengali perpindahan. Kekeliruan itu tidak menimbulkan pesan kesalahan,
sehingga hanya dapat ditemukan melalui pengukuran seperti pada praktikum ini.
\end{catatabox}
